\section{Honesty in LLMs}
\label{sec:honesty}

In the general context of human society, honesty is considered a fundamental aspect of moral character, encompassing virtues such as integrity and straightforwardness, and the absence of deceit (Dictionary, 1989). In human interactions, honesty plays a crucial role in fostering genuine connections by allowing individuals to communicate openly, which promotes interpersonal bonds, such as deeper understanding and empathy (Vivekananda \& Meenakshi, 2024), and trust (Lacey et al., 2018). In contrast, dishonest behaviors can significantly damage interpersonal relationships, as those who engage in deceit may struggle to perceive the emotions of others (Lee et al., 2019). Moreover, dishonesty can have a social contagion effect, increasing the likelihood of unethical behavior spreading among individuals (Wiltermuth et al., 2015).

Honesty is also essential in human-AI collaboration. Through interactions with AI, humans develop a mental model of the system (Hartson, 2012). Trust in the AI is strengthened when it demonstrates honest behaviors, such as expressing uncertainty (Mehrotra et al., 2024). Conversely, research indicates that dishonest behaviors, such as confidently presenting incorrect information, can undermine human trust (Dhuliawala et al., 2023b). This decline in trust persists over time and does not recover easily, even after prolonged periods (Dhuliawala et al., 2023b; Zhou et al., 2024).

\subsection{What is Honesty in LLMs}
\label{sec:what_honesty}

In the realm of LLMs, there has been a long-standing pursuit of developing honest models, with various definitions of honesty emerging over time. Askell et al. (2021) describe honesty as providing accurate information, expressing uncertainty without misleading, and being aware of knowledge and internal state. Kadavath et al. (2022) indicate honesty as an umbrella concept including truthfulness, calibration, self-knowledge, explainability, and non-deceptiveness. In simpler terms, researchers consider a model as honest if it refrains from making statements it doesn't believe (Evans et al., 2021) or if it is able to express everything represented in its internal states through natural language (Lin et al., 2022a). Recent studies suggest that an honest model should accurately express its knowledge and humbly acknowledge its limitations without deception or being inconsistent (Ward et al., 2024; Yang et al., 2024b; Chern et al., 2024).

To summarize, the most widely accepted definitions for an honest LLM are \emph{self-knowledge} and \emph{self-expression}. Self-knowledge involves the model being aware of its own capabilities, recognizing what it knows and what it doesn't, allowing it to acknowledge limitations or convey uncertainty when necessary. Self-expression refers to the model's ability to faithfully express its knowledge, leading to reliable outputs. In this paper, we consider an LLM to be honest if it fulfills these two widely accepted criteria: \emph{possessing both self-knowledge and self-expression}. An illustrated example is shown in Fig.~\ref{fig:honest_llm}, with detailed explanations provided below.

\begin{figure}[t]
\centering
\includegraphics[width=0.85\textwidth]{fig2_honest_llm.png}
\caption{An illustration of an honest LLM that demonstrates both self-knowledge and self-expression.}
\label{fig:honest_llm}
\end{figure}

\subsection{Self-knowledge}
\label{sec:self_knowledge}

The concept of self-knowledge is crucial in both the philosophy of mind and epistemology, referring to one's understanding of their own mental states, such as experience, thoughts, beliefs, and desires (Gertler, 2010). Within the context of LLMs, research on self-knowledge has also emerged as a prominent and rapidly growing field of interest (Askell et al., 2021; Kadavath et al., 2022; Yang et al., 2024b; Chern et al., 2024). Specifically, the self-knowledge capacity of LLMs hinges on their ability to recognize what they know and what they don't know. This enables them to explicitly state \emph{``I don't know''} when lacking necessary knowledge, thereby avoiding making wrong statements (Yang et al., 2024b; Zhang et al., 2024a). Additionally, it also allows them to provide confidence or uncertainty indicators in responses to reflect the likelihood of their correctness (Lin et al., 2022a; Xu et al., 2024b; Stengel-Eskin et al., 2024).

Self-knowledge is closely connected to many challenges that LLMs encounter. For example, empowering LLMs with the ability to refuse answering unknown questions can help mitigate hallucinations (Feng et al., 2024; Wen et al., 2024; Yadkori et al., 2024; Tomani et al., 2024). In addition, an LLM's uncertainty can serve as a valuable indicator for detecting hallucinations (Xiao \& Wang, 2021; Manakul et al., 2023; Kossen et al., 2024; Chen et al., 2024a; Farquhar et al., 2024). Beyond addressing hallucinations, a model's uncertainty or confidence also plays a vital role in decision-making. For instance, it can determine when external knowledge is needed in adaptive retrieval augmentation (Jiang et al., 2023; Wang et al., 2023c; Liu et al., 2024b; Ni et al., 2024; Yao et al., 2024), or whether it is necessary to invoke another LLM in a model cascading scenario (Yue et al., 2024; Jung et al., 2024; Gupta et al., 2024; Ram\'irez et al., 2024).

\subsection{Self-expression}
\label{sec:self_expression}

In human society, self-expression involves conveying one's thoughts and feelings through languages, decisions, or actions, and it is regarded as a highly respected value in Western civilization (Kim \& Ko, 2011). In the realm of LLMs, research on self-expression has gained significant attention, particularly due to the contrast between the vast knowledge acquired during the pre-training phase and the frequent occurrence of undesirable behaviors such as hallucinations (Li et al., 2024a; Lin et al., 2024a; Tian et al., 2024; Zhang et al., 2025). Specifically, we refer to self-expression as the model's ability to express its own knowledge faithfully, either based on parametric knowledge acquired through training or in-context knowledge. This enables the model to ground its responses in its knowledge rather than fabricating information.

Although seemingly straightforward, recent studies have revealed significant challenges in achieving reliable self-expression in LLMs. For example, Zhu \& Li (2023); Golovneva et al. (2024) suggest that specific data augmentation techniques during pre-training, such as paraphrasing, shuffling, or reversing, are essential for ensuring reliable knowledge expression, i.e., providing correct answers to related questions, regardless of the fine-tuning methods applied afterward. Additionally, Gekhman et al. (2024); Lin et al. (2024a); Wan et al. (2024a) demonstrate that inappropriate fine-tuning, such as introducing new knowledge in the fine-tuning dataset, can inadvertently lead to undesirable fabrications. Beyond the scope of training-time strategies, Li et al. (2024a); Chuang et al. (2024) highlight that even when an LLM possesses the knowledge internally to answer the question, frequently used decoding methods, such as greedy search, may fail to generate the correct answer. In addition to challenges with internal knowledge, LLMs often struggle to effectively convey in-context knowledge, such as information from textual documents (Liu et al., 2024e; Shi et al., 2024b; Tang et al., 2024), as well as multimodal data, including visual and audio inputs (Li et al., 2023b; Bai et al., 2024; Leng et al., 2024a), which limits their applicability.

Another reflection of the imperfect self-expression capabilities of LLMs is the frequent exhibition of inconsistent behaviors. For instance, slight changes in prompts, such as format adjustment or query paraphrasing, can affect the model's ability to express knowledge, leading to significant performance differences even with semantically equivalent prompts (Gu et al., 2023; Mizrahi et al., 2023; Wang et al., 2023b; Sclar et al., 2024; Cao et al., 2024). Additionally, LLMs may provide answers that align with the use's preferred views, even if those answers are incorrect (Perez et al., 2023; Wang et al., 2023a; Sharma et al., 2024; Huang et al., 2024b), potentially due to rewarding hacking, such as exploiting human preferences (Wei et al., 2023), when aligning the model with human values (Bai et al., 2022; Ouyang et al., 2022).
